% Generated draft: requires mathematical and source review.
\documentclass[11pt]{article}
\usepackage[margin=25mm]{geometry}
\usepackage{fontspec}
\setmainfont{Latin Modern Roman}
\usepackage{amsmath,amssymb,mathtools,unicode-math}
\setmathfont{Latin Modern Math}
\usepackage{enumitem,needspace,xurl,hyperref}
\setlength{\parindent}{0pt}
\setlength{\parskip}{5pt}
\setlength{\emergencystretch}{3em}
\newcommand{\N}{\mathbb{N}}
\newcommand{\Z}{\mathbb{Z}}
\newcommand{\Q}{\mathbb{Q}}
\newcommand{\R}{\mathbb{R}}
\newcommand{\C}{\mathbb{C}}
\newcommand{\F}{\mathbb{F}}
\newcommand{\A}{\mathbb{A}}
\newcommand{\PP}{\mathbb{P}}
\newcommand{\E}{\mathbb{E}}
\newcommand{\eps}{\varepsilon}
\newcommand{\dd}{\,\mathrm{d}}
\begin{document}
\section*{UnsolvedMath ID: 30003030}
\textbf{Source problem code:} catalogue entry
\subsection{Definitions and mathematical statement}
Let \(G\) be a finite abelian group, written additively, and let \(S=\{0=g_0,g_1,\ldots,g_{n-1}\}\subseteq G\) consist of \(n\geq 1\) distinct elements. No assumption that \(S\) generates \(G\) is imposed. Give \(\R^n\) its standard Euclidean inner product, and define
\[
A_{n-1}=\left\{(x_0,\ldots,x_{n-1})\in\Z^n:\sum_{j=0}^{n-1}x_j=0\right\},\qquad
L_G(S)=\left\{x\in A_{n-1}:\sum_{j=1}^{n-1}x_jg_j=0\text{ in }G\right\}.
\]
Here integer multiples of elements of \(G\) are interpreted using its additive group structure. Write \(V=\operatorname{span}_{\R}L_G(S)\). The defining group homomorphism has finite image, so \(L_G(S)\) has rank \(n-1\) and \(V=\{x\in\R^n:\sum_jx_j=0\}\).

For \(n\geq2\), define the minimum length, squared minimum, and set of minimal vectors by
\[
\lambda_1(L)=\min_{0\ne v\in L}\|v\|_2,\qquad
m(L)=\lambda_1(L)^2,\qquad
\operatorname{Min}(L)=\{v\in L:\|v\|_2=\lambda_1(L)\},
\quad L=L_G(S).
\]
The count \(|\operatorname{Min}(L)|\) counts opposite vectors separately. The lattice is well-rounded if \(\operatorname{span}_{\R}\operatorname{Min}(L)=V\). A basis of minimal vectors means a \(\Z\)-basis of \(L\) whose members all belong to \(\operatorname{Min}(L)\).

For a \(\Z\)-basis \(b_1,\ldots,b_{n-1}\), define the Euclidean covolume and Gram determinant by
\[
\operatorname{covol}_V(L)=\sqrt{\det\bigl(\langle b_i,b_j\rangle\bigr)_{i,j=1}^{n-1}},\qquad
\operatorname{disc}(L)=\det\bigl(\langle b_i,b_j\rangle\bigr)_{i,j=1}^{n-1}.
\]
The source does not specify whether its term ``determinant'' denotes the covolume or its square; both quantities are displayed to retain this unresolved convention. Likewise, both minimum length and squared minimum are specified because the source does not fix its norm terminology.

The full Euclidean lattice automorphism group is
\[
\operatorname{Aut}_{\mathrm{Euc}}(L)=\{T\in O(V):T(L)=L\},
\]
where \(O(V)\) is the group of linear isometries of \(V\). Its coordinate-permutation subgroup consists of restrictions to \(V\) of permutations of the \(n\) ambient coordinates that preserve \(L\). The source's intended meaning of ``automorphism group'' is not explicitly defined; the full Euclidean group and the coordinate-permutation subgroup must therefore be distinguished.

\textbf{Problem.} For every finite abelian group \(G\) and every subset \(S\) as above, determine, in terms of \(G\) and \(S\), the minimum length or squared minimum of \(L_G(S)\), its determinant under the intended convention, its number of minimal vectors, necessary and sufficient conditions for well-roundedness and for the existence of a basis of minimal vectors, and its automorphism group under the intended interpretation. This is a determination and classification problem, not an assertion that every such lattice is well-rounded or admits a basis of minimal vectors.

The permitted case \(n=1\) gives \(L_G(\{0\})=\{0\}\). There is then no nonzero vector over which to take a minimum. The source supplies no conventions for minimum-related questions or well-roundedness in rank zero; these boundary questions remain explicitly unresolved.

For \(S=G\), write \(L_G=L_G(G)\). The supplied dated assessment describes determinant, minimal-vector counts, well-roundedness, minimal bases, and a permutation part of the automorphism group for these full-group lattices as known. It describes results for general \(L_G(S)\) as limited to special families and subgroup results. These are supplied literature qualifications, not independently verified mathematical conclusions, and they do not specify the individual theorems, exceptional cases, or families involved.
\subsection{Short English statement}
For each finite abelian group \(G\) and subset \(S\) containing zero, study the integer vectors whose coordinates sum to zero and whose corresponding group-valued sum vanishes. Determine the resulting lattice's shortest vectors, determinant, spanning and basis properties, and symmetries. The source labels the problem partially solved; determinant and automorphism conventions, rank-zero behavior, and the precise scope of the reported known results require review.
\subsection{Sources}
\begin{itemize}
\item Supplied local registry record: permanent UnsolvedMath ID \(30003030\), title ``Geometry of Lattices from Finite Abelian Groups,'' source problem code \texttt{catalogue entry}, in the Oberwolfach Reports Open Problems collection. This draft uses the supplied complete record; repository files were not independently inspected.
\item Supplied source citation: \emph{Lattices and Applications in Number Theory} (2016), Oberwolfach Reports, pp.~87--154, DOI \texttt{10.4171/owr/2016/3}. Source URL and supplied literature-source link, labelled ``Oberwolfach Report 3/2016'': \url{https://doi.org/10.4171/owr/2016/3}. The record attributes its provenance to workshop reports and their open-problem discussions; no individual proposer is supplied. The proposed year is recorded as 2016.
\item Literature assessment checked on \texttt{2026-08-21}: the supplied status is ``partially solved.'' The supplied assessment reports known geometric results and a permutation part of the automorphism group for full-group lattices \(L_G\), with only special-family and subgroup results for general \(L_G(S)\). This is supplied dated curation derived from the OWR record, not independently verified here. The record expressly requests independent verification of time-sensitive status and cited claims.
\item The supplied statement-status label is \texttt{corrected\_verified}. Its accompanying qualification says that only safe notation, typography, and TeX-delimiter cleanup was applied, with all hypotheses, quantifiers, and subquestions preserved. This label does not establish independent mathematical verification of this draft.
\end{itemize}
\end{document}
